
%%%%%%%%%%%%%%%%
%%%%%%%%%%%%%%%%
%%%%%%%%%%%%%%%%
% isometric embedding  --  property (1)
%%%%%%%%%%%%%%%%
%%%%%%%%%%%%%%%%
%%%%%%%%%%%%%%%%

\subsection{Isometric embedding} \label{sec:embed}

Map each $a\in\Q$ to the class of the constant sequence $(a,a,a,\dots)$:
\[
\iota:\Q\to\Qb,\qquad \iota(a):=[(a,a,a,\dots)].
\]
\begin{align*}
&\iota(a)=\iota(b) \imp d\text{-}\lim_{n\to\infty}|a-b|=0 \imp |a-b|=0 \imp a=b
&&\text{(injective)}\\
&\iota(a)+\iota(b)=\iota(a+b);\qquad a<b\iff\iota(a)<\iota(b)
&&\text{(Step 3a, 3b)}\\
&\db(\iota(a),\iota(b))=[(d(a,b),d(a,b),\dots)]=\iota(d(a,b))
&&\text{(isometric)}
\end{align*}
Hence $\iota$ is an isometric embedding $(\Q,d)\subseteq(\Qb,\db)$, and
property (1) holds. From now on we freely regard $\Q$ as a subset of $\Qb$,
writing $a$ for $\iota(a)$ when no confusion is possible.

Distances in $\Qb$ are elements of $\Qb$, so an ``$\epsilon$'' in $\Qb$ is a
class, not a number of $\Q$. The next corollary says that this never matters:
it is enough to use $\epsilon$'s coming from $\Q$.

\begin{corollary} (wiseness) $\mid$ \label{cor:wise}
For $Y_n,Y\in\Qb$,
\[
\forall\epsilon\in\Q,\ \epsilon>0:\ \ \db(Y_n,Y)<\iota(\epsilon)\ev
\quad\iff\quad
\forall\epsilon\in\Qb,\ \epsilon>0:\ \ \db(Y_n,Y)<\epsilon\ev .
\]
We call the left side \emph{$\Q$-wise} and the right side \emph{$\Qb$-wise}
convergence; the same holds for the Cauchy condition. In particular
$\iota(e_n)\to0$ in $\Qb$.
\end{corollary}
% \begin{proof}
% \begin{align*}
% &\Leftarrow:\quad \epsilon\in\Q,\ \epsilon>0 \by{3b} \iota(\epsilon)\in\Qb,\ \iota(\epsilon)>0 .\\
% &\Rightarrow:\quad \epsilon\in\Qb,\ \epsilon>0 \by{\ref{lem:ratbounds}a}
% \exists q\in\Q:\ 0<\iota(q)<\epsilon \imp \db(Y_n,Y)<\iota(q)<\epsilon\ev .\\
% &\phantom{\Rightarrow:}\quad e_n<\epsilon\ev \by{3b}
% \db(\iota(e_n),0)=\iota(e_n)<\iota(\epsilon)\ev . \qedhere
% \end{align*}
% \end{proof}




%%% ADD cauchy subsequence lemma. 